\chapter{Welpen \& Junghunde}
\label{chap:gruppe:welpen}

\section{Sozialisation \& Spiel}

% ── Übung: Hundebegegnung ─────────────────────────────────────────────────────
\begin{uebung}{Hundebegegnung}{welpen}{Begrüßung}{2}
  \uebungsbeschreibung{%
    Zwei Welpen werden kontrolliert aneinander herangeführt. Ziel ist eine ruhige,
    entspannte Begrüßung ohne Überwältigung.\\[6pt]
    \textbf{Schritte:}
    \begin{enumerate}[leftmargin=*, itemsep=2pt]
      \item Beide Hunde an lockerer Leine, Abstand ca. 5\,m.
      \item Beide Halter gehen parallel – kein direktes Aufeinanderzugehen.
      \item Abstand schrittweise verkürzen, wenn beide Hunde entspannt sind.
      \item Kurze Begrüßung erlauben (max. 3 Sek.), dann trennen und loben.
    \end{enumerate}
  }
  \uebungsvariationen{%
    \begin{itemize}[leftmargin=*, itemsep=4pt]
      \item \textbf{Mehrere Hunde:} Übung in der Gruppe mit wechselnden Partnerhunden.
      \item \textbf{Mit Ablenkung:} Begegnung in belebter Umgebung üben.
    \end{itemize}
  }
  \uebungstrainer{%
    \begin{itemize}[leftmargin=*, itemsep=4pt]
      \item Körpersprache beider Hunde genau beobachten – bei Stress sofort Abstand vergrößern.
      \item Halter anweisen, die Leine locker zu halten – Spannung überträgt sich auf den Hund.
      \item Überwältigende Begrüßungen sofort unterbrechen.
    \end{itemize}
  }
\end{uebung}

% ── Weitere Übungen hier einfügen ─────────────────────────────────────────────
